\documentclass{article}

\usepackage[T1]{fontenc}

\usepackage{booktabs,graphicx,siunitx}

\begin{document}

\section{Random effects}

\begin{tabular}{ll}
\toprule
 & (1) \\
\midrule
(Intercept) & 0.13 \\
 & (0.47) \\
 & [$-$0.79, 1.06] \\
investment & 0.40 \\
 & (0.10) \\
 & [0.20, 0.59] \\
education\_spend & $-$1.57 \\
 & (0.24) \\
 & [$-$2.04, $-$1.10] \\
urbanization & 0.11 \\
 & (0.05) \\
 & [0.00, 0.22] \\
\midrule
Num.Obs. & 100 \\
N entities & 10 \\
R\textsuperscript{2} & 0.93 \\
R\textsuperscript{2} Adj. & 0.93 \\
Within R\textsuperscript{2} & .36 \\
Between R\textsuperscript{2} & 1.00 \\
Overall R\textsuperscript{2} & .93 \\
\bottomrule
\end{tabular}

\includegraphics{figures/01_random-effects/figure_coefficients.png}

Like fixed effects but treats entity differences as random, allowing time-invariant predictors. Read each B/p as usual - but only trust this model if the Hausman test favours random over fixed effects. Method: R plm package (Croissant \& Millo, 2008). Your significance threshold ($\alpha$) is 0.05.

In a random-effects model, predictor investment gave B=0.40, p $<$ .001.

Statistical basis: Random-effects (GLS) panel estimator (Croissant, Y., Millo, G. (2008). "Panel Data Econometrics in R: The plm Package." Journal of Statistical Software, 27(2), 1-43.).

\end{document}